\chapter{ПРОГРАММНАЯ РЕАЛИЗАЦИЯ И СРАВНИТЕЛЬНЫЙ АНАЛИЗ АЛГОРИТМОВ}
\label{ch:analysis}

Перед переходом к анализу результатов необходимо пояснить номенклатуру исследуемых алгоритмов. В задачах работы заявлено пять методов: DMD, DMDc, SINDy, SINDy-c и EDMD. Подход с добавлением индекса «c» (with control) математически является расширением базового алгоритма для систем с внешним управляющим воздействием. В разработанном программном комплексе алгоритмы реализованы в виде универсальных абстрактных классов-оберток. Если на вход алгоритму подается система без явного управления (например, автономный аттрактор Лоренца), матрица управления обнуляется, и алгоритм автоматически редуцируется до своей базовой версии (DMD, SINDy). Таким образом, на представленных ниже графиках и в таблицах названия «DMDc», «SINDy» и «EDMD» выступают в качестве имен программных модулей, аппаратно покрывающих весь заявленный спектр алгоритмов в зависимости от подаваемой на вход динамической системы.

В рамках данного раздела представлены результаты сравнительного анализа реализованных алгоритмов в унифицированной среде \texttt{dynsys-methods-evaluator}. Оценка проводилась на трех синтетических наборах данных (LTI, осциллятор Ван-дер-Поля, аттрактор Лоренца) по метрике среднеквадратичной ошибки (MSE) прогнозирования траектории. Для хаотической системы Лоренца горизонт прогнозирования был ограничен 50 шагами во избежание экспоненциального расхождения, обусловленного положительным старшим показателем Ляпунова.

\section{Влияние объема обучающей выборки на точность аппроксимации}
\label{sec:sample_size}

Объем обучающей выборки $N$ является критическим параметром для data-driven методов. Для оценки этого фактора уровень аддитивного шума был зафиксирован на отметке $\sigma = 0.05$, а объем выборки варьировался в диапазоне $N \in \{500, 1000, 2000\}$. 

\begin{figure}[ht]
    \centering
    \includegraphics[width=\textwidth]{sample_size.png}
    \caption{Влияние объема обучающей выборки на точность аппроксимации (при уровне шума $\sigma = 0.05$)}
    \label{fig:sample_size}
\end{figure}

Как видно из Рисунка \ref{fig:sample_size}, на линейной системе (LTI) метод DMDc демонстрирует стабильную работу независимо от размера выборки (MSE порядка $10^{-3}$), что подтверждает его оптимальность для линейных задач. Однако на хаотической системе Лоренца метод EDMD показал экстремальную чувствительность к недостатку данных: при $N=500$ ошибка MSE составила $68.5$, тогда как увеличение выборки до $N=2000$ позволило снизить ошибку до $0.97$. Это связано с тем, что задача линейной регрессии в бесконечномерном поднятом пространстве оператора Купмана (lifting) требует плотного заполнения фазового пространства обучающими примерами. 

Метод SINDy продемонстрировал наибольшую стабильность: при увеличении выборки с 500 до 2000 точек на системе Лоренца его ошибка упала с $0.86$ до $0.06$. Алгоритм STLSQ эффективно отсекает незначимые коэффициенты, если статистической информации достаточно для их подавления.

\section{Оценка робастности методов к уровню измерительного шума}
\label{sec:noise_robustness}

В реальных физических системах данные всегда подвержены шуму. Для анализа робастности объем данных был зафиксирован ($N=2000$), а уровень шума варьировался $\sigma \in [0.0, 0.2]$.

\begin{figure}[ht]
    \centering
    \includegraphics[width=\textwidth]{noise_robustness.png}
    \caption{Оценка робастности методов к уровню измерительного шума (при объеме выборки $N = 2000$)}
    \label{fig:noise_robustness}
\end{figure}

Графики на Рисунке \ref{fig:noise_robustness} иллюстрируют поведение алгоритмов при зашумленных наблюдениях. Линейный метод DMDc не смог аппроксимировать нелинейные системы (Ван-дер-Поль, Лоренц). Алгоритм расходился, возвращая некорректные предсказания, что подтверждает его математическую неприменимость для систем со сложной топологией без использования ядерных расширений.

Метод EDMD показал склонность к многократному усилению шума. На системе Ван-дер-Поля при $\sigma=0$ ошибка составляла $0.31$, но при $\sigma=0.2$ она возросла до $2.58$. Использование полиномиального словаря наблюдаемых приводит к тому, что гауссовский шум возводится в квадрат, формируя тяжелые выбросы, с которыми не справляется $L_2$-регуляризация.

Метод SINDy оказался наиболее робастным. На системе Лоренца при максимальном шуме ($\sigma=0.2$) MSE составила всего $0.12$. Заложенное в метод предположение о разреженности динамики работает как мощный фильтр шума, позволяя выделить истинные управляющие уравнения на фоне искажений наблюдений.

\section{Статистический анализ результатов и границы применимости}
\label{sec:stats_conclusion}

Проведенное исследование позволяет четко очертить эмпирические границы применимости рассмотренных методов:
\begin{enumerate}
    \item Метод DMD/DMDc является оптимальным выбором для линейных систем, обеспечивая скорость и точность, однако он строго неприменим к хаотическим и нелинейным осцилляторам.
    \item Метод EDMD способен описывать нелинейную динамику, однако его использование ограничено низким уровнем измерительного шума. В условиях сильных зашумлений полиномиальное расширение фазового пространства разрушает предиктивную способность оператора Купмана.
    \item Метод SINDy является наиболее универсальным подходом, обеспечивающим высокую точность даже на хаотических аттракторах при высоком шуме, при условии наличия репрезентативной выборки ($N \ge 2000$) для работы алгоритма.
\end{enumerate}

Таким образом, для прикладных задач, где измерения сопряжены с существенным шумом сенсоров, приоритетным подходом следует считать алгоритм SINDy.